\begin{proofbox}
	\textbf{\textcolor{CProof}{思路提示}}
	利用函数图像的点集定义，通过分析图像上点的坐标变换来导出解析式的关系。

	\medskip
	\textbf{\textcolor{CProof}{正式推导}}
	\begin{description}[style=nextline, leftmargin=2.8em, labelsep=0.5em, itemsep=0.55em]
		\item[\textbf{步骤一（取点列式）：}]
			在原函数$y=f(x)$的图像上任取一点$P(x_0,y_0)$，则$y_0=f(x_0)$。
			\par\textbf{理由：} 函数图像上的点一定满足函数解析式。
		\item[\textbf{步骤二（向上平移）：}]
			将点$P$向上平移$b$个单位，得到点$P'(x_0,y_0+b)$。
			\par\textbf{理由：} 向上移动$b$个单位使纵坐标增加$b$，横坐标保持不变。
		\item[\textbf{步骤三（建立关系）：}]
			点$P'$在平移后的函数图像上，因此对任意$x_0$，平移后函数$g$满足$g(x_0)=y_0+b=f(x_0)+b$，从而$g(x)=f(x)+b$。
			\par\textbf{理由：} 平移后的图像由所有这样的点构成，纵坐标即为新函数在对应$x$处的值。
		\item[\textbf{步骤四（向下同理）：}]
			类似地，将点$P$向下平移$b$个单位，可得新函数为$g(x)=f(x)-b$。
			\par\textbf{理由：} 向下平移使纵坐标减少$b$，其余推导完全相同。
	\end{description}

	\medskip
	\textbf{\textcolor{CProof}{结论回扣}}
	因此，函数$y=f(x)$的图像沿$y$轴向上平移$b$个单位后的解析式为$y=f(x)+b$；向下平移$b$个单位后的解析式为$y=f(x)-b$。这就是“上加下减”的数学本质。
\end{proofbox}
